\subsubsection{FORMATIONS}

\begin{enumerate}
  \item The Formation a unit or stack of units is in plays a very important role in the game. In Column Formation, units are deployed for rapid marching and manuevering, but are not efficient for combat. In Battle Formation, units are deployed for maximum effectiveness in combat, but are not efficient for rapid movement.

  \item Changing the formation a unit counter is in from Column Formation to Battle Formation or back again costs a varying number of Movement Allowance Factors, depending on the type of unit and the terrain in the hex in which the Formation change takes place. These costs are shown on the MOVEMENT \& TERRAIN EFFECTS CHART.

  \item In hexes where there are several types of terrain features, the highest cost is the one used.

  \textit{For example, a unit in a hex containing both clear and slope terrain will use up the number of Movement Allowance Factors appropriate for slope terrain when changing Formation.}

  \item Roads are ignored as terrain for the process of changing Formation.

  \item A player may change the Formation his units are in at any time during his Turn, provided that the units have enough remaining Movement Allowance Factors to do so.

  \item It is possible for some units in a hex to be in Column Formation, while other units in the same hex are in Battle Formation. Note that such a situation will cause some varying movement costs through the same hex. Even though in different Formations, all units in the same hex must be stacked together.

  It is possible for the Phasing Player to change Formation while in an enemy Zone of Control. This must be done prior to beginning the first Round of Combat.
\end{enumerate}

\subsubsection{MARKING FORMATIONS}

\begin{enumerate}
  \item Units in Column Foramtion are indicated by placing a Column Formation Marker on top of a unit or Stack of units in that Formation.

  \item Units in Battle Formation are indicated by \textit{not} having a Column Formation marker placed on top of them, i.e. any units that do not have a Column Formation Marker placed on top of them are in Battle Formation.

  \item For hexes containing units in both Formations, indicate by placing Column Formation Marker on top of them. Then place the units in Battle Formation at the top of the Stack, on top of the Column Formation Marker.
\end{enumerate}